\section{A project that can explain itself}

\subsection{The project-understanding Assistant}

The project-understanding Assistant is an authenticated C-LARA-2 feature
through which users can ask unconstrained natural-language questions about the
platform itself. It is neither a conventional FAQ system nor a manually
maintained database of answers. Instead, Codex receives the question in the
read-only restricted environment described in Section~8, inspects the current
repository, and returns an answer grounded in source code, tests, roadmaps,
issue records and operational documentation. References to repository files
and line numbers are converted into clickable links, allowing a user to
inspect the evidence behind an answer.

The deployment constraints described in Section~8 are central to this design.
The Assistant needs enough access to inspect a current repository checkout,
but it must not receive general authority to modify the repository or control
the production server. The resulting arrangement makes repository-grounded
explanation available inside the platform while constraining the powers of the
agent providing it.

\subsection{Repository-grounded answers and inspectable evidence}

Different repository artefacts answer different parts of a project question.
Source code and tests describe current implemented behaviour; roadmaps record
the history and intended direction of larger areas of work; issue records
track focused problems and unresolved tasks; and how-to files describe
operations that human collaborators may need to perform. The Assistant can
therefore address not only whether a capability is implemented, but also why a
design was chosen, what remains planned, and what workaround is currently
available.

This capability depends on the shared project memory described in Section~7.
The Assistant does not acquire information by introspection or make a claim to
human-like comprehension. In a deliberately limited operational sense, it can
inspect the artefacts through which the project records its implementation,
history and planned development, and construct a context-sensitive,
evidence-linked account for users and collaborators.

\subsection{Implemented capability: Mandarin and pinyin}

A first example, shown in Figure~\ref{fig:mandarin-pinyin}, concerns whether
the platform can create Mandarin texts and add pinyin annotations. The
Assistant answers that Mandarin is supported as a
project and text-pipeline language, then locates evidence across several parts
of the repository: language choices and project models, the Mandarin
text-generation prompt, tokenisation using \texttt{jieba}, pinyin processing
using \texttt{pypinyin}, manual editing support, and ruby-style HTML
rendering. This is a useful test because no single file contains the complete
answer; the Assistant must connect several implementation layers. The
annotation architecture underlying those layers is described in Section~4.

\begin{figure}[htbp]
\centering
\begin{minipage}{\linewidth}
\centering
\includegraphics[width=\linewidth]{images/assistant_mandarin_pinyin_question.jpg}\par\medskip
\includegraphics[width=\linewidth]{images/assistant_mandarin_pinyin_answer.jpg}
\end{minipage}
\caption{A project-understanding Assistant interaction concerning Mandarin
text generation and pinyin annotation. The answer is constructed by inspecting
several parts of the current repository and includes links to the supporting
source files.}
\label{fig:mandarin-pinyin}
\end{figure}

\subsection{Honest answers about planned work}

A second example, shown in Figure~\ref{fig:haitian-audio}, concerns
native-speaker audio recording for Haitian Creole, where current synthetic
speech may be unsuitable. The Assistant distinguishes
what is implemented from what is only planned. It identifies the present
no-audio/skip-TTS option, explains that recording through the platform is not
yet an implemented feature, and links to the relevant roadmap and issue
record for the planned community recording/upload workflow. This kind of
answer is valuable because it gives a useful, evidence-linked ``not yet''
response rather than presenting a roadmap aspiration as current functionality.

\begin{figure}[htbp]
\centering
\begin{minipage}{\linewidth}
\centering
\includegraphics[width=\linewidth]{images/assistant_haitian_audio_question.jpg}\par\medskip
\includegraphics[width=\linewidth]{images/assistant_haitian_audio_answer.jpg}
\end{minipage}
\caption{An Assistant response to a request for native-speaker audio recording
in Haitian Creole. The answer separates current functionality, planned work
and a partial workaround, with links to source code, roadmaps and issue
records.}
\label{fig:haitian-audio}
\end{figure}

\subsection{Multilingual questions}

Although much early testing used English questions, the Assistant can also
respond usefully to questions in other major languages. Figure~\ref{fig:french-roadmap-issues}
shows a French question asking why C-LARA-2 uses roadmap and issue files in
the repository: \emph{Pourquoi C-LARA-2 utilise-t-il des fichiers de feuille
de route et des fichiers d'incidents dans le dépôt ?} In English, this asks:
``Why does C-LARA-2 use roadmap files and issue files in the repository?''

The answer remains in French and explains that roadmaps provide structured
planning and longer-term direction: objectives, architecture, phases,
priorities and a record of what is complete, in progress or planned. Issue
records provide canonical operational tracking of focused problems and tasks,
including state, priority, dependencies and current work. It then explains
why the separation is useful: humans steer priorities while Codex maintains an
auditable repository memory linking longer-term planning to an executable,
versioned short-term backlog. The answer links to the repository evidence that
supports this account. This example reinforces the point that the Assistant is
not limited to a fixed English help system; it can construct repository-grounded
explanations in the language of the question.

\begin{figure}[htbp]
\centering
\begin{minipage}{\linewidth}
\centering
\includegraphics[width=\linewidth]{images/assistant_french_roadmap_issues_question.jpg}\par\medskip
\includegraphics[width=\linewidth]{images/assistant_french_roadmap_issues_answer.jpg}
\end{minipage}
\caption{A French project-understanding question and answer about the
complementary roles of roadmaps and issue files. The Assistant answers in
French and links to the relevant repository evidence.}
\label{fig:french-roadmap-issues}
\end{figure}

\subsection{Explaining the report itself}

Once the current report had been committed to the repository and deployed with
the current platform version, the Assistant could also answer a question about
the report itself. Figure~\ref{fig:report-speech-recognition} shows a user who
had been skimming the report asking whether C-LARA-2 can use speech recognition
and, if so, how to invoke it. The response correctly answers that speech
recognition is not an implemented C-LARA-2 feature. It links to Section~10,
which states that spoken learner input is not currently supported; to
Section~11, which presents speech recognition for whole-text alignment as
future work; and to the current audio pipeline, which offers text-to-speech or
no-audio/skip-TTS modes but no ASR/transcription stage.

This is a useful self-referential test. The answer does not merely retrieve a
sentence from the report: it reconciles the report's statements about planned
work with source-code evidence about the current implementation, then gives a
clear answer to the practical question of what a user can invoke now. It also
illustrates why a repository-resident report can become part of the project's
usable memory rather than remaining a static publication artefact.

\begin{figure}[htbp]
\centering
\begin{minipage}{\linewidth}
\centering
\includegraphics[width=\linewidth]{images/assistant_report_speech_recognition_question.jpg}\par\medskip
\includegraphics[width=\linewidth]{images/assistant_report_speech_recognition_answer.jpg}
\end{minipage}
\caption{A self-referential Assistant interaction concerning the Initial
C-LARA-2 Project Report. The answer distinguishes speech recognition as future
work from the currently implemented TTS and no-audio modes, using both the
report and current repository code as evidence.}
\label{fig:report-speech-recognition}
\end{figure}

\subsection{Limits, review and future evidence records}

The Assistant cannot be perfect. Documentation may lag behind code, preserve
an abandoned intention, or omit a relevant relationship; a cited file may be
relevant without being sufficient; and an answer can be accurate without being
optimally helpful. Its answers should therefore be treated as informed,
inspectable accounts rather than definitive statements. Important answers will
need a future workflow for export, versioning and human assessment if they are
to become durable evidence for project claims.

Nevertheless, when code, tests, documentation, roadmaps and issues are
accessible to the same constrained agent, the repository becomes more than a
container for software. Its external memory becomes actively usable: the
Assistant can help human collaborators and users navigate current behaviour,
history, plans and limitations, while making its evidence available for
inspection.
